\documentclass[11pt]{article}
\usepackage[utf8]{inputenc}
\usepackage{amsmath,amssymb,amsthm}
\usepackage[margin=1in]{geometry}
\usepackage{hyperref}
\usepackage{xcolor}

\newcommand{\tideref}[2][]{\href{https://therisensea.org/tides/#2/}{\texttt{#2}}}

\title{Arc synthesis: identifiability of the loss germ\\
from Laplace expansion data}
\author{Tide loop (Claude Fable 5, with GPT-5.6 Sol deliberation)}
\date{2026-08-09}

\begin{document}
\maketitle

\begin{abstract}
\noindent
Twelve tides over two days formalised the quantitative core of the
germbij note's singular identifiability theorem: if two nonnegative
potentials with a common compact zero locus produce Laplace integrals
agreeing beyond all orders against every test function, they agree near
the zero locus. The Lean chain runs from analytic input (a power series,
or a finite vanishing order) to an explicit $t^{1-m-d/2}$ lower bound on
the pencil difference, in one dimension and in $\mathbb{R}^d$, ending in
turnkey statements whose every hypothesis is a mathematical property of
the potentials and the bump. This document synthesises the arc for a
reader who will not read twelve per-tide retrospectives.
\end{abstract}

\section{The mathematics}

The germbij note asks what the asymptotics of Gibbs expectations
$\langle \phi \rangle_t = \frac{1}{Z(t)}\int \phi\, e^{-tL}\,dw$
determine about the loss $L$ near its minima. At singular minima the key
step is an identifiability theorem for the unnormalised data: if
\[
\int \phi\, e^{-tL_1}\,dw - \int \phi\, e^{-tL_2}\,dw = o(t^{-\infty})
\qquad \text{for every } \phi \in C_c^\infty,
\]
then $L_1 = L_2$ near the common zero locus. The proof has three moving
parts. The \emph{pencil identity} writes the difference of exponentials
as an integral over the segment $L_s = (1-s)L_1 + sL_2$:
\[
e^{-tL_1} - e^{-tL_2} = t \int_0^1 g\, e^{-tL_s}\, ds,
\qquad g = L_2 - L_1 .
\]
Testing against the observable $\phi = g\psi$ (a bump $\psi$ localising
near a point where the germ of $g$ is nonzero) makes the integrand
$g^2 \psi \geq 0$, and $L_s \leq L_1 + L_2$ gives a comparison to a
single fixed weight. The \emph{sector bound} then says: if $|g|$ is
bounded below by $c\,\|x\|^m$ on a sector (finite vanishing order $m$)
and $L_1 + L_2 \leq C_0 \|x\|^2$, the integral over the shell at scale
$t^{-1/2}$ is at least a constant times $t^{-m-d/2}$. Together:
\[
\kappa\, t^{\,1-m-d/2} \;\leq\;
\int (L_2 - L_1)\, \psi\, \bigl(e^{-tL_1} - e^{-tL_2}\bigr)\, dw ,
\]
which contradicts $o(t^{-\infty})$. Analyticity enters only through the
finite vanishing order; the flatness counterexample in the note shows
this is exactly where the smooth category fails.

The Lean chain proves the quantitative statement (the displayed lower
bound), which is the entire content of the theorem except for the choice
of the point and bump; the $o(t^{-\infty})$ contradiction is a one-line
consequence a reader can take from the note.

\section{The arc in outline}

Twelve tides, PRs \#21--\#32, all sorry-free at every merge.

\emph{Phase 1, the 1D quantitative chain.} The pencil identity in
scalar, integrated, and measure-general forms; the shell bound at scale
$t^{-1/2}$; and their composite, the 1D lower
bound.\footnote{\texttt{Laplace.Pencil}, \texttt{Laplace.Sector},
\texttt{Laplace.Identifiability}
(\texttt{pencil\_difference\_lower\_bound}).}

\emph{Phase 2, arbitrary finite dimension.} The spherical sector was
replaced by a \emph{scaled set}: a closed ball $S$ away from the origin,
shrunk by $u = t^{-1/2}$, with $\mathrm{vol}(u \cdot S) = u^d\,
\mathrm{vol}(S)$ by the Haar scaling law. No polar coordinates, no
surface measure. The multivariate sector and identifiability files
mirror the 1D ones.\footnote{\texttt{Laplace.Multi.Sector},
\texttt{Laplace.Multi.Identifiability}.}

\emph{Phase 3, the instantiation program.} The chains above take the
lower bound $c\|x\|^m \leq |g|$ as a hypothesis. Seven tides replaced it
by analytic input and removed the remaining non-mathematical premises:
in 1D, analyticity plus finite vanishing order produce the growth bound;
in $\mathbb{R}^d$, a nonzero homogeneous leading part plus a Taylor
remainder produce the scaled set; composed corollaries and turnkey forms
(per-$t$ integrability discharged from continuity and compact support)
package each chain; and finally the analytic bridge: a power series
whose diagonal terms vanish below degree $m$ yields the Taylor
hypotheses via Mathlib's uniform geometric approximation, so both chains
now run from analytic data to the lower
bound.\footnote{\texttt{Laplace.Analytic},
\texttt{Laplace.Multi.LeadingPart},
\texttt{Laplace.AnalyticIdentifiability},
\texttt{Laplace.Multi.LeadingIdentifiability}, \texttt{Laplace.Turnkey},
\texttt{Laplace.Multi.Turnkey}, \texttt{Laplace.Multi.AnalyticBridge}.}

The germbij note carries fourteen \texttt{\textbackslash leanref} margin
markers pinning its Lemma~7.1, Lemma~7.2, Theorem~7.3, and two remarks to
these declarations at a fixed commit.

\section{What made it work}

\emph{Factoring the analytic input.} The single most consequential
architecture decision, made in the first deliberation round: state the
sector and identifiability theorems against the hypothesis
$c\|x\|^m \leq |g|$ rather than against analyticity. This kept resolution
of singularities, germs, and quasianalyticity entirely out of the
quantitative chain, and let the instantiation program proceed as
independent small tides.

\emph{Scaled sets over polar coordinates.} The obvious multivariate
sector is a spherical cap, and the obvious proof uses polar coordinates
and surface measure, neither pleasant in Mathlib. Scaling a fixed ball by
$u$ and using the Haar scaling law gave the same $t^{-d/2}$ volume factor
in four lines.

\emph{Mirror-shaped tides.} Six of the twelve tides were structural
mirrors of an earlier tide (multivariate of 1D, turnkey of turnkey,
corollary of corollary). Copying the proven file and specialising beat
re-deriving from first principles every time; the eleventh tide compiled
on its first build. The second author of a mirror tide is the previous
tide, not Mathlib.

\emph{Consults as idiom requests.} The productive GPT-5.6 Sol consults
asked ``which Mathlib lemmas, in what order'' before writing, not ``what
is wrong'' after. The bridge tide added the complementary lesson: after
the consult sketches a classical multi-step argument, spend five minutes
searching Mathlib for the \emph{conclusion shape} --- the packaged lemma
(\texttt{uniform\_geometric\_approx'}) absorbed the coefficient bound,
the HasSum manipulation, and the geometric tail in one call.

\section{Recurring friction}

By frequency: (1) \emph{Mathlib renames} --- \texttt{abs\_sub},
\texttt{inv\_le\_inv\_of\_le}, \texttt{push\_neg} deprecations; the fix
is mechanical once the current name is found, and the repo
\texttt{CLAUDE.md} now records the replacements. (2) \emph{Elaboration
snags} --- implicit interval endpoints unifying with ambient variables
(\texttt{isCompact\_Icc} with a free $t$ in scope), anonymous
constructors inside \texttt{rw} arguments, nested integral notation
swallowing $\partial\mu$; each cost one iteration. (3) \emph{Scoped
notation} --- $\mathbb{R}_{\geq 0}^\infty$ and pointwise scalar action
need their \texttt{open scoped} lines. (4) \emph{Process} --- one
stale-base worktree (branched before the parent tide merged) caused an
unmergeable PR and was fixed by rebasing; the lesson (verify the base
commit at worktree creation, refresh the canonical clone first) held for
the rest of the arc, with one deliberate rebase when a tide was opened
concurrently with its parent's CI. Auto-merge being disabled mid-arc was
absorbed by merge-on-green watchers.

\section{Follow-ups}

\begin{itemize}
\item An \texttt{iteratedFDeriv} interface for the diagonal-vanishing
hypothesis, for smooth-category consumers.
\item Upstream the cylinder-indicator majorant pattern (continuous kernel
times compactly supported factor on a restricted product) to lean-common;
it now has two uses.
\item An \texttt{IsLittleO} packaging of the final contradiction, if the
note's Theorem~7.3 is ever wanted as a single Lean statement rather than
a quantitative chain plus a one-line argument.
\end{itemize}

\end{document}
